\documentclass[12pt,a4paper]{article}
\usepackage{mathwizards-ingles}
\title{%
  \vspace{-1cm}
  {\Huge\bfseries\textcolor{mes3}{Mes 3 · Semana 12}}\\[0.3cm]
  {\Large Second Conditional, \textit{Used to} y Chunks Situacionales}\\[0.2cm]
  {\large\color{grissuave} Nivel B1 · Consolidación Final e Independencia}\\[0.5cm]
  \rule{\textwidth}{1.5pt}
}
\author{}
\date{}

\begin{document}
\maketitle
\thispagestyle{empty}

\begin{cajanivel}
\textbf{Objetivo:} Semana final del programa. El estudiante consolida su nivel B1 con el \textbf{Second Conditional} (situaciones hipotéticas), la estructura \textbf{\textit{used to / would}} para hábitos pasados, \textbf{conectores avanzados} y \textbf{chunks situacionales} para role-play de situaciones reales (entrevistas, viajes, restaurantes, emergencias).
\end{cajanivel}

% ═══════════════════════════════════════════════════════════════════════════════
\section{Gramática Emergente: Second Conditional}
% ═══════════════════════════════════════════════════════════════════════════════

\begin{cajagrammar}[Second Conditional: situaciones hipotéticas / irreales]
Se usa para hablar de situaciones \textbf{imaginarias o improbables} en el presente o futuro.

\medskip
\textbf{Estructura:} \eng{If} + Past Simple, \eng{would} + verbo base.

\begin{tabularx}{\textwidth}{@{} X @{}}
\toprule
\textbf{Ejemplos} \\
\midrule
\textit{\textbf{If} I \textbf{had} a million dollars, I \textbf{would travel} the world.} \\
\textit{\textbf{If} I \textbf{were} you, I \textbf{would study} more.} \\
\textit{\textbf{If} she \textbf{spoke} English fluently, she \textbf{would get} a better job.} \\
\textit{I \textbf{would buy} a house \textbf{if} I \textbf{had} enough money.} \\
\bottomrule
\end{tabularx}

\medskip
\begin{cajaejemplo}[Notas Importantes]
\begin{itemize}[nosep, leftmargin=0.8cm]
  \item Después de ``if'', se usa \textbf{Past Simple}, pero NO se habla del pasado: es una situación \textbf{imaginaria ahora}.
  \item Con \textit{to be}, la forma correcta para todas las personas es \textbf{``were''} (no ``was''):\\
    $\checkmark$ \textit{If I \textbf{were} rich...} \quad (formal/correcto)\\
    $\sim$ \textit{If I \textbf{was} rich...} \quad (coloquial, aceptable)
  \item Se puede usar \textbf{could} en lugar de \textit{would}:\\
    \textit{If I had more time, I \textbf{could} learn Japanese.}
\end{itemize}
\end{cajaejemplo}
\end{cajagrammar}

\begin{cajagrammar}[Contraste: First vs. Second Conditional]
\begin{tabularx}{\textwidth}{@{} l X X @{}}
\toprule
& \textbf{First Conditional} & \textbf{Second Conditional} \\
\midrule
\textbf{Probabilidad} & Real / Probable & Imaginaria / Improbable \\
\textbf{Tiempos} & If + Present, will + V & If + Past, would + V \\
\textbf{Ejemplo} & \textit{If it rains, I'll take an umbrella.} & \textit{If I were you, I'd take an umbrella.} \\
& (puede llover) & (yo no soy tú: imaginario) \\
\bottomrule
\end{tabularx}
\end{cajagrammar}

% ═══════════════════════════════════════════════════════════════════════════════
\section{Gramática Emergente: \textit{Used to} y \textit{Would}}
% ═══════════════════════════════════════════════════════════════════════════════

\begin{cajagrammar}[\textit{Used to} + verbo base: hábitos pasados que ya no existen]
\textbf{Estructura:} Sujeto + \eng{used to} + verbo base.

\medskip
\begin{tabularx}{\textwidth}{@{} l X @{}}
\toprule
\textbf{Forma} & \textbf{Ejemplo} \\
\midrule
\textbf{Afirmativo} & \textit{I \textbf{used to} play football when I was young.} \\
\textbf{Negativo} & \textit{She \textbf{didn't use to} like vegetables.} \\
\textbf{Pregunta} & \textit{\textbf{Did} you \textbf{use to} live in Bogotá?} \\
\bottomrule
\end{tabularx}

\medskip
\begin{cajaejemplo}[\textit{Would} para hábitos pasados repetidos]
\textbf{Would} también expresa hábitos pasados, pero \textbf{solo para acciones} (no para estados):

\begin{itemize}[nosep, leftmargin=0.8cm]
  \item $\checkmark$ \textit{We \textbf{would} always play in the park after school.} (acción repetida)
  \item $\times$ \textit{I \textbf{would} live in Bogotá.} (``live'' es un estado $\rightarrow$ usar \textit{used to})
  \item $\checkmark$ \textit{I \textbf{used to} live in Bogotá.}
\end{itemize}
\end{cajaejemplo}
\end{cajagrammar}

% ═══════════════════════════════════════════════════════════════════════════════
\section{Chunks: Conectores Avanzados}
% ═══════════════════════════════════════════════════════════════════════════════

\begin{cajachunk}[Advanced Connectors — Nivel B1]
\begin{tabularx}{\textwidth}{@{} l l X @{}}
\toprule
\textbf{Conector} & \textbf{Función} & \textbf{Ejemplo} \\
\midrule
\eng{despite} / \eng{in spite of} & concesión + sustantivo & \textit{\textbf{Despite} the rain, we went out.} \\
\eng{although} / \eng{even though} & concesión + cláusula & \textit{\textbf{Although} I was tired, I studied.} \\
\eng{whereas} & contraste directo & \textit{She likes coffee, \textbf{whereas} he prefers tea.} \\
\eng{as long as} & condición & \textit{You can go \textbf{as long as} you come back early.} \\
\eng{unless} & a menos que & \textit{I'll go \textbf{unless} it rains.} \\
\eng{in order to} & propósito & \textit{She studies \textbf{in order to} pass the exam.} \\
\eng{due to} & causa formal & \textit{The game was cancelled \textbf{due to} bad weather.} \\
\eng{not only ... but also} & adición enfática & \textit{He is \textbf{not only} smart \textbf{but also} kind.} \\
\bottomrule
\end{tabularx}

\medskip
\begin{cajaejemplo}[Despite vs. Although]
\textbf{despite / in spite of} + \textbf{sustantivo o -ing}: \textit{Despite \textbf{being} tired...}\\
\textbf{although / even though} + \textbf{cláusula completa}: \textit{Although I \textbf{was} tired...}
\end{cajaejemplo}
\end{cajachunk}

% ═══════════════════════════════════════════════════════════════════════════════
\section{Chunks Situacionales: Role-Play}
% ═══════════════════════════════════════════════════════════════════════════════

\begin{cajachunk}[At a Restaurant — En un Restaurante]
\begin{tabularx}{\textwidth}{@{} X l @{}}
\eng{A table for two, please.} & \esp{Una mesa para dos, por favor.} \\
\eng{Can I see the menu, please?} & \esp{¿Puedo ver el menú?} \\
\eng{I'd like to order...} & \esp{Me gustaría pedir...} \\
\eng{What do you recommend?} & \esp{¿Qué recomiendas?} \\
\eng{Could I have the bill, please?} & \esp{¿Me trae la cuenta?} \\
\eng{Is service included?} & \esp{¿Está incluido el servicio?} \\
\eng{Can I pay by card?} & \esp{¿Puedo pagar con tarjeta?} \\
\eng{The food was delicious, thank you.} & \esp{La comida estuvo deliciosa, gracias.} \\
\end{tabularx}
\end{cajachunk}

\begin{cajachunk}[At a Job Interview — En una Entrevista de Trabajo]
\begin{tabularx}{\textwidth}{@{} X l @{}}
\eng{I'm here for the interview.} & \esp{Vengo a la entrevista.} \\
\eng{I have experience in ...} & \esp{Tengo experiencia en ...} \\
\eng{My strengths are ...} & \esp{Mis fortalezas son ...} \\
\eng{I used to work as a ...} & \esp{Trabajé como ...} \\
\eng{I'm a team player.} & \esp{Trabajo bien en equipo.} \\
\eng{I'm looking for an opportunity to grow.} & \esp{Busco una oportunidad para crecer.} \\
\eng{When would I start?} & \esp{¿Cuándo empezaría?} \\
\eng{Thank you for your time.} & \esp{Gracias por su tiempo.} \\
\end{tabularx}
\end{cajachunk}

\begin{cajachunk}[Traveling — Viajando]
\begin{tabularx}{\textwidth}{@{} X l @{}}
\eng{I'd like to book a room.} & \esp{Quisiera reservar una habitación.} \\
\eng{I have a reservation under...} & \esp{Tengo una reserva a nombre de...} \\
\eng{What time is check-in / check-out?} & \esp{¿A qué hora es la entrada/salida?} \\
\eng{Where is the nearest pharmacy?} & \esp{¿Dónde queda la farmacia más cercana?} \\
\eng{I'd like a window seat, please.} & \esp{Quisiera un asiento de ventana.} \\
\eng{My flight has been delayed.} & \esp{Mi vuelo ha sido retrasado.} \\
\eng{Where can I exchange money?} & \esp{¿Dónde puedo cambiar dinero?} \\
\end{tabularx}
\end{cajachunk}

\begin{cajachunk}[Emergencies — Emergencias]
\begin{tabularx}{\textwidth}{@{} X l @{}}
\eng{Help! / Please help me!} & \esp{¡Ayuda! / ¡Por favor, ayúdenme!} \\
\eng{Call an ambulance!} & \esp{¡Llamen una ambulancia!} \\
\eng{I need a doctor.} & \esp{Necesito un médico.} \\
\eng{I've lost my passport.} & \esp{He perdido mi pasaporte.} \\
\eng{My wallet was stolen.} & \esp{Me robaron la cartera.} \\
\eng{Where is the hospital?} & \esp{¿Dónde está el hospital?} \\
\eng{I'm allergic to ...} & \esp{Soy alérgico/a a ...} \\
\eng{It's an emergency.} & \esp{Es una emergencia.} \\
\end{tabularx}
\end{cajachunk}

% ═══════════════════════════════════════════════════════════════════════════════
\section{Oraciones Listas para Minar}
% ═══════════════════════════════════════════════════════════════════════════════

\begin{cajamineria}[Oraciones \textit{i+1} — Semana 12]

\begin{enumerate}[leftmargin=1.2cm]
  \item \textit{If I \underline{were} you, I would accept the offer.}\\
    {\small\textbf{Objetivo:} Second Conditional con ``if I were you'' (consejo hipotético).}

  \item \textit{If I \underline{had} more time, I would learn to play the guitar.}\\
    {\small\textbf{Objetivo:} Second Conditional para deseos irreales.}

  \item \textit{I \underline{used to} live in a small town, but now I live in the city.}\\
    {\small\textbf{Objetivo:} ``used to'' para hábitos/estados pasados.}

  \item \textit{When we were kids, we \underline{would} always play in the garden.}\\
    {\small\textbf{Objetivo:} ``would'' para acciones repetidas en el pasado.}

  \item \textit{\underline{Despite} being very tired, she finished the marathon.}\\
    {\small\textbf{Objetivo:} ``despite'' + -ing para concesión.}

  \item \textit{He passed the exam \underline{even though} he didn't study much.}\\
    {\small\textbf{Objetivo:} ``even though'' + cláusula completa.}

  \item \textit{I'll go to the party \underline{unless} I have to work late.}\\
    {\small\textbf{Objetivo:} ``unless'' como ``if not'' condicional.}

  \item \textit{Could I have the \underline{bill}, please? Everything was excellent.}\\
    {\small\textbf{Objetivo:} chunk de restaurante: ``Could I have the bill''.}

  \item \textit{I \underline{used to} be afraid of flying, but now I love it.}\\
    {\small\textbf{Objetivo:} ``used to'' con un adjetivo/estado emocional.}

  \item \textit{She is \underline{not only} a great singer \underline{but also} a talented writer.}\\
    {\small\textbf{Objetivo:} la estructura enfática ``not only ... but also''.}
\end{enumerate}
\end{cajamineria}

% ═══════════════════════════════════════════════════════════════════════════════
\section*{Nota Final: ¿Qué Sigue Después del B1?}
% ═══════════════════════════════════════════════════════════════════════════════

\begin{cajaconsolidacion}[Guía para Continuar tu Inmersión (B2+)]
Completar el nivel B1 significa que eres un \textbf{usuario independiente}. Para seguir avanzando:

\begin{enumerate}[leftmargin=1.2cm]
  \item \textbf{Consumo nativo:} Transicionar de material adaptado a \textbf{contenido nativo} (series, películas, libros sin adaptar, podcasts de interés general).
  \item \textbf{Producción oral:} Buscar compañeros de conversación o tutores en línea (iTalki, HelloTalk).
  \item \textbf{Escritura:} Comenzar a escribir diarios, correos o ensayos cortos en inglés.
  \item \textbf{Anki:} Mantener repasos diarios, pero la proporción de estudio deliberado baja al 10--15\%. El 85--90\% debe ser inmersión pura.
  \item \textbf{Gramática avanzada:} Third Conditional, Mixed Conditionals, Subjunctive, inversiones, cleft sentences — todo esto se adquiere mejor a través de lectura extensiva que de estudio explícito.
\end{enumerate}
\end{cajaconsolidacion}

\vfill
\begin{center}
\rule{0.5\textwidth}{0.5pt}\\[0.3cm]
{\small\color{grissuave} Curso de Inmersión en Inglés · Mes 3 · Semana 12 · Nivel B1 · Consolidación Final}
\end{center}

\end{document}
